% Appendix B

\chapter{Draft Paper Outline}
\label{AppendixB}
\lhead{Appendix B. \emph{Draft Paper Outline}}

\section{Draft Title}
\textbf{IronClaw: Dynamic Tool Synthesis with Secure MicroVM Isolation for
Autonomous Agents}

\section{Draft Abstract}
Autonomous agents based on Large Language Models increasingly rely on external
tools to complete long-horizon software engineering, data processing, and
operational tasks. Fixed tool registries simplify control and auditing, but they
limit the agent when a task requires an unanticipated parser, transformation,
test harness, or integration. Dynamic tool synthesis addresses this limitation
by allowing the agent to generate executable code during a task, but doing so
introduces a severe security problem: generated code is untrusted and may be
incorrect, malicious, prompt-injected, or resource-exhaustive.

This paper will present IronClaw, a self-hosted agent platform that separates
trusted host control from untrusted guest execution. The host daemon manages
model interaction, policy, memory, configuration, and sandbox lifecycle, while a
guest runtime executes generated tools inside a restricted environment. The
system supports process, container, and Firecracker-backed microVM runtimes over
a shared protobuf protocol. The planned evaluation will compare fixed-tool and
dynamic-tool agents on multi-step tasks, measure containment against malicious
payloads, and quantify latency and resource overhead across sandbox backends.
The expected contribution is a practical architecture and empirical methodology
for making runtime tool synthesis useful without collapsing the host trust
boundary.

\section{Planned Paper Sections}
\begin{enumerate}
    \item \textbf{Introduction}
    \begin{itemize}
        \item Motivation for dynamic tool synthesis.
        \item Security risk of executing agent-generated code.
        \item Summary of IronClaw's host/guest architecture.
        \item Contributions and evaluation goals.
    \end{itemize}

    \item \textbf{Background and Related Work}
    \begin{itemize}
        \item LLM agents and tool use.
        \item Dynamic skill and tool acquisition.
        \item Sandboxing, containers, virtual machines, and microVMs.
        \item Firecracker and lightweight virtualization.
    \end{itemize}

    \item \textbf{Threat Model}
    \begin{itemize}
        \item Untrusted generated code.
        \item Prompt injection and malicious task data.
        \item Host filesystem, network, memory, and credential protection.
        \item Non-goals such as proving the absence of all hypervisor bugs.
    \end{itemize}

    \item \textbf{System Design}
    \begin{itemize}
        \item Host daemon responsibilities.
        \item Guest runtime responsibilities.
        \item Protobuf message envelope and capability handshake.
        \item Tool allowlists and policy gates.
        \item Runtime options: process, container, and microVM.
    \end{itemize}

    \item \textbf{Implementation}
    \begin{itemize}
        \item Rust workspace organization.
        \item Configuration model.
        \item Transport and framing.
        \item Firecracker manager and vsock setup.
        \item Tool registry, file tools, restricted shell, browser tool, memory,
        and scheduler.
    \end{itemize}

    \item \textbf{Evaluation Methodology}
    \begin{itemize}
        \item Utility benchmarks for fixed-tool versus dynamic-tool agents.
        \item Security red-team payload suite.
        \item Performance measurements for latency, CPU, memory, and startup.
        \item Reproducibility and environment details.
    \end{itemize}

    \item \textbf{Results}
    \begin{itemize}
        \item Task completion rate.
        \item Turn count and token usage.
        \item Containment outcomes.
        \item Runtime overhead comparisons.
    \end{itemize}

    \item \textbf{Discussion}
    \begin{itemize}
        \item When dynamic synthesis helps.
        \item When fixed tools are sufficient.
        \item Security and usability trade-offs.
        \item Operational limits of microVM-based isolation.
    \end{itemize}

    \item \textbf{Limitations}
    \begin{itemize}
        \item Benchmark representativeness.
        \item Hardware dependency.
        \item Model dependency.
        \item Scope of red-team payloads.
    \end{itemize}

    \item \textbf{Conclusion}
    \begin{itemize}
        \item Summary of architecture, measurements, and lessons.
        \item Future work on stronger policies, snapshots, caching, and broader
        benchmarks.
    \end{itemize}
\end{enumerate}
